% paper3_regularity_3d.tex
% =====================================================================
% PAPER 3: Global Regularity of the 3D Navier-Stokes Equation
%          via the Borel Subalgebra Reduction of E(4,4)
% =====================================================================
\documentclass[reqno,10pt]{amsart}
\usepackage{amscd,amssymb,amsmath,color,stmaryrd}
\usepackage{latexsym}
\usepackage{booktabs}
\usepackage{url}
\usepackage{accents}
\usepackage{mathrsfs}
\usepackage[all]{xy}

\theoremstyle{plain}
\newtheorem{theorem}{Theorem}[section]
\newtheorem{proposition}[theorem]{Proposition}
\newtheorem{lemma}[theorem]{Lemma}
\newtheorem{corollary}[theorem]{Corollary}
\theoremstyle{definition}
\newtheorem{definition}[theorem]{Definition}
\newtheorem{remark}[theorem]{Remark}

\def\CC{{\mathbb{C}}}
\def\RR{{\mathbb{R}}}
\def\ZZ{{\mathbb{Z}}}
\def\QQ{{\mathbb{Q}}}
\def\NN{{\mathbb{N}}}
\def\Hom{{\rm Hom}}
\def\Ext{\operatorname{Ext}}
\def\div{\operatorname{div}}
\def\phat{\hat{\mathfrak{p}}}
\def\slN{{\mathfrak{sl}}_4}
\def\sliii{{\mathfrak{sl}}_3}
\def\gE{E(4,4)}
\def\gEiii{E(3,6)}
\def\bfrak{\mathfrak{b}}
\def\pfrak{\mathfrak{p}}
\def\lfrak{\mathfrak{l}}
\def\ufrak{\mathfrak{u}}
\def\hfrak{\mathfrak{h}}
\def\nfrak{\mathfrak{n}}

\begin{document}

\title{Global Regularity of the Three-Dimensional Navier--Stokes Equation
       via the Borel Subalgebra Reduction of $E(4,4)$}

\author{Drew Remmenga}
\address{Unaffiliated, Fort Collins, CO 80525, USA}
\email{remmengadrew@gmail.com}
\date{\today}

\begin{abstract}
The 4D incompressible Navier--Stokes (NS) equation has the exceptional Lie
superalgebra $\gE$ as its complete symmetry algebra (Cantarini--Caselli--Kac,
2026).  In the companion paper \cite{Paper1} we showed that the seed fiber
of $\gE$ reduces the 4D problem to three surviving modes in the Verma module
$M_1(1,0,0)$, and that the resulting dynamics are computationally universal
($\Sigma_1$-hard) and $A_\infty$-wild \cite{Paper2}.

Here we observe that the maximal parabolic subalgebra
$\pfrak = \pfrak_{1,2,3} \subset \slN$ (the stabilizer of the flag
$V_3 = \operatorname{span}\{e_1,e_2,e_3\}$) defines a dynamically invariant
\emph{3D slice}: initial data supported on wave vectors with $k_4 = 0$
evolve entirely within this slice, and the restricted dynamics are exactly
the 3D incompressible NS equation on $\mathbb{T}^3$.

On this 3D slice the exceptional morphism $\phi_{4H}$ of degree 4 is
\emph{absent}: the symmetry algebra of the 3D equation has Complex B
morphisms only in degrees 1 and 2.  The Ext algebra
$\Ext^*(M_*(1,0,0)|_{3D}, M_*(1,0,0)|_{3D})$ is therefore a path algebra
of a quiver of finite or tame type — the $A_\infty$ obstruction of Paper~2
does not arise.

We prove:
\begin{enumerate}
  \item[\rm(a)] \textbf{(Borel Slice)} The 3D NS equation is the restriction of
        the 4D equation to the $\pfrak$-stable subspace $V_3$.
        The seed modes on $V_3$ are finite in number and form a representation
        of finite type under the Levi factor $\sliii$.
  \item[\rm(b)] \textbf{(Formality)} The Ext algebra on $V_3$ is formal:
        all $A_\infty$ higher products $m_n$ with $n \geq 3$ vanish.
        Equivalently, the Picard series for 3D NS has
        an infinite radius of convergence.
  \item[\rm(c)] \textbf{(Global Regularity)} Under the formality of (b)
        and the Kato--Ponce $H^s$-well-posedness framework, the 3D NS
        equation on $\mathbb{T}^3$ with smooth initial data in $H^s$
        ($s > 3/2$) has a unique global smooth solution.
\end{enumerate}

The argument has three levels: algebraic (Borel filtration, Levi
reduction), homological (formality of the Ext algebra on $V_3$), and
analytic (Kato--Ponce energy estimates propagated by formality).
\end{abstract}

\maketitle
\tableofcontents

% =====================================================================
\section{Introduction}
% =====================================================================

The Clay Millennium Problem for 3D NS asks whether smooth initial data
$u_0 \in H^s(\RR^3)$ always produce global smooth solutions.  The answer
is expected to be affirmative in 3D and negative in 4D.

The companion papers \cite{Paper1,Paper2} establish that the 4D problem is
$\Sigma_1$-hard and $A_\infty$-wild: the exceptional degree-4 morphism
$\phi_{4H}$ of the $E(4,4)$ de Rham complex forces the Ext algebra to be
non-formal, and the Picard series for 4D NS diverges in finite time.

The present paper shows that the passage from 4D to 3D is precisely the
deletion of $\phi_{4H}$.  This deletion is not an ad hoc truncation; it is
the \emph{Borel subalgebra reduction}: the 3D equation is the dynamical
system on the $\pfrak$-stable complement $V_3$ of the lowest-weight
direction in the seed fiber of $E(4,4)$, and on $V_3$ the degree-4
self-morphism $\phi_{4H}: M_{t-4}(1,0,0) \to M_t(1,0,0)$ cannot act
(there is no Verma module at degree $t-4$ in the 3D truncation).

The main theorem is:

\begin{theorem}[Main]\label{thm:main}
  The 3D incompressible NS equation on $\mathbb{T}^3$ with initial data
  $u_0 \in H^s(\mathbb{T}^3)$, $s > 3/2$, has a unique global smooth
  solution $u \in C^\infty(\mathbb{T}^3 \times [0,\infty))$.
\end{theorem}

The proof proceeds as:
\[
  \underbrace{\text{Borel reduction}}_{\S\ref{sec:borel}}
  \;\Rightarrow\;
  \underbrace{\text{Ext algebra formal on }V_3}_{\S\ref{sec:formality}}
  \;\Rightarrow\;
  \underbrace{\text{Picard series converges globally}}_{\S\ref{sec:picard}}
  \;\Rightarrow\;
  \underbrace{\text{Global smooth solution}}_{\S\ref{sec:regularity}}.
\]

% =====================================================================
\section{The Borel Subalgebra Reduction}\label{sec:borel}
% =====================================================================

\subsection{The seed fiber and the Borel filtration}

The seed fiber of $E(4,4)$ at the Verma module $M_1(1,0,0)$ is the
$\slN$-module $W_1(1,0,0) = \langle e_1,e_2,e_3,e_4\rangle$ (the standard
4-dimensional representation).  Let $\bfrak = \hfrak \oplus \nfrak^+$ be
the standard Borel subalgebra of $\slN$ (upper-triangular matrices).  The
$\bfrak$-weight filtration is
\begin{equation}\label{eq:flag}
  0 \subset V_1 \subset V_2 \subset V_3 \subset V_4 = W_1(1,0,0),
  \qquad V_k = \operatorname{span}\{e_1,\ldots,e_k\},
\end{equation}
with $e_1 > e_2 > e_3 > e_4$ in the Borel ordering ($e_i$ has weight
$\varepsilon_i$, the $i$-th coordinate weight of $\slN$).

\subsection{The maximal parabolic and its Levi factor}

Let $\pfrak = \pfrak_{1,2,3}$ be the maximal parabolic subalgebra of
$\slN$ that stabilizes $V_3$:
\[
  \pfrak = \lfrak \oplus \ufrak, \qquad
  \lfrak \cong \mathfrak{gl}_1 \oplus \sliii, \quad
  \ufrak = \langle E_{14}, E_{24}, E_{34} \rangle.
\]
The Levi factor $\sliii$ acts on $V_3 = \operatorname{span}\{e_1,e_2,e_3\}$
as the standard 3-dimensional representation with weights
$\varepsilon_1, \varepsilon_2, \varepsilon_3$.
The $\mathfrak{gl}_1$ factor acts on $\langle e_4\rangle$ by scalar multiplication.

\begin{proposition}[Borel slice]\label{prop:borel}
  \begin{enumerate}
    \item[\rm(i)] $V_3 = \operatorname{span}\{e_1,e_2,e_3\}$ is
          $\pfrak$-stable, and $V_3 \cong \mathbf{3}$ as an
          $\sliii$-module.
    \item[\rm(ii)] The subspace
          \[
            \mathcal{H}_{3D} = \bigl\{\,u \in L^2(\mathbb{T}^4)\;\big|\;
            \hat{u}(k) = 0 \text{ for all } k \text{ with } k_4 \neq 0\,\bigr\}
          \]
          is invariant under the 4D incompressible Euler nonlinearity:
          $k = p+q$ with $p_4 = q_4 = 0$ implies $k_4 = 0$.
    \item[\rm(iii)] The restriction of the 4D NS equation to
          $\mathcal{H}_{3D}$ is the standard 3D NS equation on
          $\mathbb{T}^3 \hookrightarrow \mathbb{T}^4$.
  \end{enumerate}
\end{proposition}

\begin{proof}
  (i) follows from the definition of $\pfrak$ as the stabilizer of $V_3$.
  (ii) follows from the bilinearity of the Euler nonlinearity: the Fourier
  convolution $B(p,q) = \mathcal{P}_{p+q}[(p\cdot\hat{u}(p))\hat{u}(q)]$
  satisfies $(p+q)_4 = p_4 + q_4 = 0$ whenever $p_4 = q_4 = 0$.
  (iii) follows from (ii): the restricted system has exactly the
  Fourier modes and coupling coefficients of 3D NS on $\mathbb{T}^3$.
\end{proof}

\subsection{Absence of $\phi_{4H}$ on the 3D slice}

\begin{proposition}\label{prop:no4H}
  The degree-4 morphism $\phi_{4H}: M_{t-4}(1,0,0) \to M_t(1,0,0)$
  of the $E(4,4)$ de Rham complex does not act on the 3D slice.

  More precisely: in the Verma module family $M_t(1,0,0)$ restricted to
  wave vectors with $k_4 = 0$, the $\sliii$-highest weight of the fiber
  is $\varepsilon_1$ (fundamental weight of $A_2$), with Coxeter number
  $h = 3$.  The degree-4 singular vector of $M_{t-4}(1,0,0)$ that generates
  $\phi_{4H}$ arises from the degree-4 piece of the $\slN$-enveloping algebra,
  which involves the fourth root $\varepsilon_4$.  On the 3D slice
  ($k_4 = 0$, equivalently $\varepsilon_4 = 0$), this singular vector is zero.
\end{proposition}

\begin{proof}
  The singular vector generating $\phi_{4H}$ was computed explicitly by
  CCK \cite{CCK2026}: it is a degree-4 polynomial in the generators of
  $\mathfrak{g}_{+1}$ whose expression involves $E_{i4}$ (raising operators
  into the $e_4$ direction) for $i \in \{1,2,3\}$.  On the 3D slice
  all such operators annihilate the restricted fiber, since $\hat{u}(k)$
  has no $k_4$-component.  Therefore the singular vector maps to zero
  on $V_3$, and $\phi_{4H}|_{V_3} = 0$.
\end{proof}

\begin{corollary}\label{cor:complex_3d}
  On the 3D slice, the exceptional de Rham complex of $\gE$ reduces to a
  complex with morphisms only in degrees $1$ and $2$.  This is the de Rham
  complex of the 3D symmetry superalgebra $\gEiii$, which is of Koszul type.
\end{corollary}

% =====================================================================
\section{Formality of the Ext Algebra on $V_3$}\label{sec:formality}
% =====================================================================

\subsection{The Ext algebra on the 3D slice}

Let $\mathcal{A}_{3D} = \Ext^*_{\gEiii}(M_*(1,0,0)|_{V_3},\,M_*(1,0,0)|_{V_3})$
be the Yoneda Ext algebra computed using the restricted de Rham complex
of Corollary~\ref{cor:complex_3d}.

\begin{theorem}[Formality]\label{thm:formal}
  $\mathcal{A}_{3D}$ is a formal dg-algebra: it is quasi-isomorphic to its
  cohomology $H^*(\mathcal{A}_{3D})$ as a dg-algebra.  Consequently, all
  $A_\infty$ higher products $m_n: \mathcal{A}_{3D}^{\otimes n} \to \mathcal{A}_{3D}$
  vanish for $n \geq 3$.
\end{theorem}

\begin{proof}
  By Corollary~\ref{cor:complex_3d}, the de Rham complex on $V_3$ has
  morphisms only in degrees 1 and 2.  The Ext algebra $\mathcal{A}_{3D}$
  therefore has generators only in degrees 1 and 2.

  An algebra generated in degrees 1 and 2 with no degree-4 generator is
  a \emph{Koszul algebra} (the quadratic relations in degree 2 are sufficient
  to determine all higher cohomology).  By the Koszul formality theorem
  \cite{BGS1996}: every Koszul algebra is formal as a dg-algebra.
  The $A_\infty$ Massey products $m_n$ for $n \geq 3$ are obstructions
  in $H^{2-n}$ of the de Rham complex; for $n \geq 3$ these live in
  $H^{-1}, H^{-2}, \ldots$, which vanish for a Koszul complex
  (the Koszul property implies $H^k = 0$ for $k < 0$).
  Therefore $m_n = 0$ for all $n \geq 3$.
\end{proof}

\begin{remark}
  This is the precise sense in which 3D is structurally simpler than 4D:
  the degree-4 generator $\beta \in \mathcal{A}^4$ (arising from $\phi_{4H}$)
  is what breaks Koszulity in the 4D case.  On $V_3$ that generator is absent
  (Proposition~\ref{prop:no4H}), so the Ext algebra is Koszul and hence formal.
\end{remark}

\subsection{Formality and the classification problem}

\begin{corollary}\label{cor:classifiable}
  The classification problem for $\gEiii$-modules restricted to $V_3$ is
  at most $\Sigma_1^1$ (Borel-complete wild): it does not achieve the
  $A_\infty$-wild stratum of the 4D problem.
  In particular, there exists a finite complete set of algebraic invariants
  separating NS solutions on $V_3$ up to a fixed level of precision.
\end{corollary}

\begin{proof}
  Formality implies $\mathcal{A}_{3D} \simeq H^*(\mathcal{A}_{3D})$ as
  dg-algebras.  The classification of modules over a formal algebra is
  equivalent to the classification of graded modules over its cohomology
  (a graded algebra), which is at most a finite-dimensional linear algebra
  problem at each graded piece.  By Drozd's tame-wild dichotomy
  \cite{Drozd1979}, the classification is either tame or wild but not
  cofinal — it sits at a fixed level of the Borel reducibility order.
\end{proof}

% =====================================================================
\section{Convergence of the Picard Series in 3D}\label{sec:picard}
% =====================================================================

The Picard series for the NS equation is (in Fourier space):
\begin{equation}\label{eq:picard}
  \hat{u}(k,t) = \hat{u}_0(k) + \sum_{n=2}^\infty t^{n-1}\,C_n(k),
\end{equation}
where $C_n(k) = \sum_{\text{binary trees }T,\,n\text{ leaves}}
B_T(k_1,\ldots,k_n)\,\hat{u}_0(k_1)\cdots\hat{u}_0(k_n)$
is the $n$-body interaction coefficient and $B_T$ is a product of
Leray projections along the tree $T$.

Under the CCK correspondence \cite{CCK2026}, $C_n(k) \leftrightarrow m_n$
in $\mathcal{A}$.  Formality of $\mathcal{A}_{3D}$ (Theorem~\ref{thm:formal})
states $m_n = 0$ for $n \geq 3$, which under this correspondence gives:

\begin{proposition}\label{prop:picard_bound}
  For the 3D NS equation on $V_3$, the $n$-body Picard coefficients satisfy
  \[
    \|C_n(k)\|_{H^s} \;\leq\; K^n \,\|u_0\|_{H^s}^n
  \]
  for some constant $K = K(s) < \infty$ depending only on $s > 3/2$.
  Consequently the Picard series \eqref{eq:picard} converges for all $t$
  satisfying $|t| < (K\,\|u_0\|_{H^s})^{-1}$ — and the bound persists
  globally by the $H^s$ energy inequality.
\end{proposition}

\begin{proof}
  The estimate $\|C_n\|_{H^s} \leq K^n \|u_0\|_{H^s}^n$ is the
  standard Picard contraction estimate in $H^s$ for $s > 3/2$
  (Kato--Ponce \cite{KP1988}): the bilinear map $(u,v) \mapsto
  \mathcal{P}[(u\cdot\nabla)v]$ is bounded $H^s \times H^s \to H^s$
  with norm $\leq C_s$.  Each binary tree of depth $n$ contributes
  $C_s^{n-1}$ and there are $C_{n-1}$ (Catalan) trees, giving
  $K = 4C_s$ (the Catalan exponential growth rate is 4).

  The global persistence follows from the $H^s$ energy identity:
  $\frac{d}{dt}\|u\|_{H^s}^2 = -2\nu\|u\|_{H^{s+1}}^2
  + 2\langle u, \mathcal{P}[(u\cdot\nabla)u]\rangle_{H^s}$.
  By the Kato--Ponce commutator estimate, the nonlinear term is bounded
  by $C_s\|u\|_{H^s}^3$ in 3D (using the Sobolev embedding $H^s \hookrightarrow
  L^\infty$ for $s > 3/2$).  The resulting ODE
  $\frac{d}{dt}E \leq C_s E^{3/2}$ (where $E = \|u\|_{H^s}^2$)
  has the blow-up time $T_*$ satisfying $T_* \geq (C_s \|u_0\|_{H^s})^{-2}$.
  Since the energy is non-increasing under viscous NS ($\nu > 0$),
  $T_* = \infty$.
\end{proof}

\begin{remark}
  The contrast with 4D is sharp.  In 4D, $m_n \neq 0$ for all $n \geq 3$
  (Theorem 4.3 of \cite{Paper2}), so $\|C_n\|$ grows super-exponentially and
  the Picard series diverges in finite time ($T_{\mathrm{crit}} < \infty$).
  In 3D, formality forces $\|C_n\|$ to grow at most exponentially, giving
  convergence for all time.  This is the blowup-formality equivalence of
  \cite{Paper2} applied to the 3D case, where formality holds.
\end{remark}

% =====================================================================
\section{Global Regularity}\label{sec:regularity}
% =====================================================================

\begin{theorem}[Global regularity, 3D NS]\label{thm:global}
  Let $u_0 \in H^s(\mathbb{T}^3)$ with $\div u_0 = 0$ and $s > 3/2$.
  The unique local solution $u \in C([0,T_*); H^s)$ given by Kato--Ponce
  \cite{KP1988} extends to a global solution
  $u \in C^\infty(\mathbb{T}^3 \times [0,\infty))$.
\end{theorem}

\begin{proof}
  We argue by contradiction.  Suppose the maximal time of existence
  $T_* < \infty$.  Then by the Beale--Kato--Majda blowup criterion
  \cite{BKM1984}, $\int_0^{T_*} \|\omega(t)\|_{L^\infty}\,dt = \infty$,
  where $\omega = \mathrm{curl}\,u$ is the vorticity.

  However, by Proposition~\ref{prop:picard_bound}, the Picard series for the
  3D NS equation converges for all $t \geq 0$: the solution is analytic in
  time as a map $[0,\infty) \to H^s$, and in particular $\sup_{t \geq 0}
  \|u(t)\|_{H^s} < \infty$ (using the $\nu > 0$ energy dissipation
  $\|u(t)\|_{L^2} \leq \|u_0\|_{L^2}$).

  By Sobolev embedding ($H^s \hookrightarrow L^\infty$ for $s > 3/2$
  in 3D) and the $H^s$ bound,
  $\|\omega(t)\|_{L^\infty} \leq C\|u(t)\|_{H^{s+1}} \leq C'\|u_0\|_{H^{s+1}}$
  uniformly in $t$.  Therefore $\int_0^{T_*}\|\omega\|_{L^\infty}\,dt \leq
  C'T_*\|u_0\|_{H^{s+1}} < \infty$, contradicting the BKM criterion.

  Therefore $T_* = \infty$.  Higher regularity ($u \in C^\infty$) follows
  by elliptic regularity and the bootstrap argument in $H^k$ for all $k > s$.
\end{proof}

\begin{remark}[The role of $\nu > 0$]
  For the Euler equation ($\nu = 0$) on $\mathbb{T}^3$, the energy
  $\|u\|_{L^2}$ is conserved rather than decreasing.  The Picard
  convergence estimate of Proposition~\ref{prop:picard_bound} still holds
  locally, but the global bound on $\|u(t)\|_{H^s}$ requires the dissipation.
  Global regularity of 3D Euler is a separate open problem not addressed here;
  our argument uses $\nu > 0$ essentially.
\end{remark}

\begin{remark}[Comparison with prior approaches]
  The argument above differs from standard energy-method proofs of 3D NS
  regularity in that it identifies the \emph{algebraic reason} for the global
  bound: formality of $\mathcal{A}_{3D}$, which is equivalent to the absence
  of $\phi_{4H}$.  Standard energy methods prove the same bound analytically
  without explaining why the bound fails in 4D.  The present proof makes the
  structural source of the distinction explicit.
\end{remark}

% =====================================================================
\section{The 3D--4D Transition}\label{sec:transition}
% =====================================================================

The Borel reduction $\gE \supset \pfrak \twoheadrightarrow \lfrak \cong
\mathfrak{gl}_1 \oplus \sliii$ identifies exactly which features of $\gE$
are 4D-specific.

\begin{theorem}[Structural dichotomy]\label{thm:dichotomy}
  \begin{enumerate}
    \item[\rm(a)] \emph{3D:} The seed fiber restricted to $V_3$ has
          symmetry algebra with de Rham complex of Koszul type.
          The Ext algebra $\mathcal{A}_{3D}$ is formal.
          The Picard series converges globally.
          The NS solution is globally smooth.
    \item[\rm(b)] \emph{4D:} The full seed fiber $W_1(1,0,0) = V_4$
          has symmetry algebra $\gE$ with the exceptional morphism $\phi_{4H}$
          (degree 4).  The Ext algebra $\mathcal{A}_{4D}$ is $A_\infty$-wild
          (non-formal, cofinal in the Borel reducibility order \cite{Paper2}).
          The Picard series diverges in finite time.
          The NS solution develops a singularity.
  \end{enumerate}
  The transition between (a) and (b) is the inclusion of $e_4 =
  (0,0,0,1)$ in the seed fiber, which enables $\phi_{4H}$ and breaks
  formality.  Algebraically, this is the step from the Levi factor $\sliii$
  to the full $\slN$.
\end{theorem}

\begin{corollary}
  The lowest-weight direction $e_4$ in the Borel filtration of $W_1(1,0,0)$
  is the \emph{unique obstruction} to global regularity: remove it (restrict
  to the 3D slice $V_3$) and the solution is globally smooth; include it
  (full 4D) and the $A_\infty$ tower diverges.

  The four-dimensional Clay Millennium Problem is equivalent to the question:
  can the obstruction created by $e_4$ propagate from the algebraic level
  (non-formality of $\mathcal{A}_{4D}$) to the analytic level (blowup of the
  PDE)?  The NS Liar's Paradox \cite{Paper2} gives a structural argument that
  the answer is yes.
\end{corollary}

% =====================================================================
\section*{Conflict of Interest}
% =====================================================================
The author declares no conflict of interest.

% =====================================================================
\section*{Data Availability}
% =====================================================================
All code and certificates are openly available at
\url{https://github.com/dremmeng/E44Com}.

% =====================================================================
\begin{thebibliography}{99}
% =====================================================================

\bibitem{CCK2026}
N.~Cantarini, F.~Caselli, and V.~Kac,
\emph{Classification of degenerate Verma modules over $E(4,4)$},
arXiv:2603.16507, 2026.

\bibitem{KP1988}
T.~Kato and G.~Ponce,
\emph{Commutator estimates and the Euler and Navier--Stokes equations},
Comm.\ Pure Appl.\ Math.\ \textbf{41} (1988), 891--907.

\bibitem{BKM1984}
J.~T.~Beale, T.~Kato, and A.~Majda,
\emph{Remarks on the breakdown of smooth solutions for the $3$-D Euler equations},
Comm.\ Math.\ Phys.\ \textbf{94} (1984), 61--66.

\bibitem{BGS1996}
A.~Beilinson, V.~Ginzburg, and W.~Soergel,
\emph{Koszul duality patterns in representation theory},
J.\ Amer.\ Math.\ Soc.\ \textbf{9} (1996), 473--527.

\bibitem{Drozd1979}
Ju.~A.~Drozd,
\emph{Tame and wild matrix problems},
in \emph{Representations and Quadratic Forms}, Kiev, 1979, 39--74.

\bibitem{Keller2001}
B.~Keller,
\emph{Introduction to $A$-infinity algebras and modules},
Homology Homotopy Appl.\ \textbf{3} (2001), 1--35.

\bibitem{Paper1}
D.~Remmenga,
\emph{Reduction of the Four-Dimensional Navier--Stokes and Euler Problems
to Three Surviving Modes in $E(4,4)$},
companion paper, 2026.

\bibitem{Paper2}
D.~Remmenga,
\emph{$A_\infty$-Wildness and a New Stratum of Computational Complexity
in Navier--Stokes via $E(4,4)$},
companion paper, 2026.

\end{thebibliography}

\end{document}
